\paragraph{有理玫瑰读出、universal solenoid 极限与虚部可见层的核消隐}
在有限素数局部化群的 prime-register solenoid、无穷 prime-register 的有限可见因子化，以及 RH 的 Carath\'eodory--Toeplitz 视界接口均已建立之后，还可把有理玫瑰族、Lissajous 族与分母扫描的几何读出压缩为下列一条连续链。其核心是：经典平面曲线可视为字符读出在圆商与奇异环之间的阴影，而相位补全后的扫描坐标则把分母层级严格识别为 universal solenoid 的有限层因子。

\begin{theorem}[有限 prime-register solenoid 上的 Kronecker 读出与有理玫瑰/Lissajous 的统一提升]\label{thm:xi-time-part62d-solenoid-kronecker-rose-lissajous-lift}
\leanverified{paper\_xi\_time\_part62d\_solenoid\_kronecker\_rose\_lissajous\_lift}
沿用定理 \ref{thm:killo-localization-circle-dimension-prime-ledger-splitting} 的记号。设 \(S\subset\mathbb P\) 为有限素数集，并记
\[
G_S:=\ZZ[S^{-1}],
\qquad
\Sigma_S:=\widehat{G_S},
\qquad
0\longrightarrow K_S:=\prod_{p\in S}\ZZ_p
\longrightarrow
\Sigma_S
\xrightarrow{\ \pi\ }
\TT
\longrightarrow 0.
\]
定义 Kronecker 读出轨道
\[
\kappa_S:\RR\longrightarrow \Sigma_S,
\qquad
\bigl(\kappa_S(t)\bigr)(q):=e^{2\pi iqt}
\qquad (q\in G_S).
\]
并对每个 \(q\in G_S\) 记字符
\[
\chi_q:\Sigma_S\to\TT,
\qquad
\chi_q(x):=x(q).
\]
则成立：
\begin{enumerate}
  \item \(\kappa_S\) 是连续单射群同态，且 \(\overline{\kappa_S(\RR)}=\Sigma_S\)。
  \item 对任意 \(q\in G_S\)，有理玫瑰轨道
  \[
  \gamma_q(\theta):=\cos(q\theta)e^{i\theta}
  \qquad (\theta\in\RR)
  \]
  满足严格恒等式
  \[
  \gamma_q(\theta)
  =
  \frac12\Bigl(
  \chi_{1+q}\bigl(\kappa_S(\theta/2\pi)\bigr)
  +
  \chi_{1-q}\bigl(\kappa_S(\theta/2\pi)\bigr)
  \Bigr).
  \]
  \item 对任意整数 \(a,b\in\ZZ\) 与相位 \(\delta\in\RR\)，Lissajous 轨道
  \[
  L_{a,b,\delta}(t):=(\cos(at+\delta),\cos(bt))
  \]
  满足
  \[
  L_{a,b,\delta}(t)
  =
  \Bigl(
  \Re\bigl(e^{i\delta}\chi_a(\kappa_S(t/2\pi))\bigr),\
  \Re\bigl(\chi_b(\kappa_S(t/2\pi))\bigr)
  \Bigr).
  \]
\end{enumerate}
\end{theorem}
\begin{proof}
把 \(\Sigma_S\) 视为 \(\TT^{G_S}\) 的闭子群，则 \(\kappa_S\) 的每个坐标函数
\[
t\longmapsto e^{2\pi iqt}
\qquad (q\in G_S)
\]
都连续，故 \(\kappa_S\) 连续；又显然
\[
\kappa_S(t+t')=\kappa_S(t)\kappa_S(t'),
\]
故其为群同态。

若 \(\kappa_S(t)=\kappa_S(t')\)，则对 \(q=1\) 有
\[
e^{2\pi i(t-t')}=1,
\]
从而 \(t-t'\in\ZZ\)。再对任意
\[
N=\prod_{p\in S}p^{k_p}
\]
取 \(q=1/N\in G_S\)，得到
\[
e^{2\pi i(t-t')/N}=1,
\]
故 \(t-t'\in N\ZZ\)。由于此结论对所有这样的 \(N\) 同时成立，只能有 \(t-t'=0\)。因此 \(\kappa_S\) 单射。

对稠密性，只需计算其湮没子。若 \(q\in G_S\) 满足
\[
\chi_q(\kappa_S(t))=1
\qquad
(\forall\, t\in\RR),
\]
则
\[
e^{2\pi iqt}=1
\qquad
(\forall\, t\in\RR),
\]
于是必有 \(q=0\)。故 \(\kappa_S(\RR)\) 的湮没子平凡。由 Pontryagin 对偶中“闭子群由其湮没子确定”的标准结论，\(\overline{\kappa_S(\RR)}=\Sigma_S\)。

对第 \(2\) 条，直接代入
\[
\chi_{1\pm q}\bigl(\kappa_S(\theta/2\pi)\bigr)=e^{i(1\pm q)\theta}
\]
即得
\[
\frac12\Bigl(e^{i(1+q)\theta}+e^{i(1-q)\theta}\Bigr)
=
e^{i\theta}\frac{e^{iq\theta}+e^{-iq\theta}}{2}
=
\cos(q\theta)e^{i\theta}.
\]

对第 \(3\) 条，同理有
\[
\chi_a(\kappa_S(t/2\pi))=e^{iat},
\qquad
\chi_b(\kappa_S(t/2\pi))=e^{ibt}.
\]
取实部便得到
\[
\Re\bigl(e^{i\delta}\chi_a(\kappa_S(t/2\pi))\bigr)=\cos(at+\delta),
\qquad
\Re\bigl(\chi_b(\kappa_S(t/2\pi))\bigr)=\cos(bt).
\]
证毕。
\end{proof}

\begin{proposition}[universal solenoid 的相位补全分母扫描与有限层因子化]\label{prop:xi-time-part62d-universal-solenoid-phase-completed-scan}
\leanverified{paper\_xi\_time\_part62d\_universal\_solenoid\_phase\_completed\_scan}
记
\[
\Sigma_{\mathbb P}:=\widehat{\QQ}.
\]
则有短正合列
\[
0\longrightarrow \prod_{p\in\mathbb P}\ZZ_p
\longrightarrow \Sigma_{\mathbb P}
\xrightarrow{\ \pi\ }
\TT
\longrightarrow 0,
\]
并且
\[
\Sigma_{\mathbb P}
\cong
\varprojlim_{N\ge 1}\widehat{(1/N)\ZZ}
\cong
\varprojlim_{N\ge 1}\TT,
\]
其中当 \(N\mid M\) 时连接映射为 \(z\mapsto z^{M/N}\)。

对任意 \(q\in\QQ\)，定义相位补全后的扫描坐标
\[
\Phi_q^{\cos}(x):=\frac12\bigl(\chi_{1+q}(x)+\chi_{1-q}(x)\bigr),
\qquad
\Phi_q^{\sin}(x):=\frac{1}{2i}\bigl(\chi_{1+q}(x)-\chi_{1-q}(x)\bigr).
\]
则成立：
\begin{enumerate}
  \item 对每个 \(q\in\QQ\)，有
  \[
  \Phi_q^{\cos}(x)+i\Phi_q^{\sin}(x)=\chi_{1+q}(x).
  \]
  \item 映射
  \[
  \Phi^{\mathrm{ph}}:\Sigma_{\mathbb P}\longrightarrow \prod_{q\in\QQ}\CC^2,
  \qquad
  x\longmapsto \bigl((\Phi_q^{\cos}(x),\Phi_q^{\sin}(x))\bigr)_{q\in\QQ}
  \]
  是拓扑嵌入。
  \item 若 \(E\subset\QQ\) 为有限集，且对每个 \(q\in E\) 记 \(d(q)\) 为其既约正分母，并令
  \[
  N_E:=\operatorname{lcm}\{d(q):\ q\in E\},
  \qquad
  H_{N_E}:=(1/N_E)\ZZ,
  \]
  则有限投影
  \[
  \Phi_E^{\mathrm{ph}}:\Sigma_{\mathbb P}\longrightarrow \prod_{q\in E}\CC^2
  \]
  仅依赖于自然商映射
  \[
  \pi_{N_E}:\Sigma_{\mathbb P}\twoheadrightarrow \widehat{H_{N_E}}\cong \TT.
  \]
\end{enumerate}
因此，任意有限分母扫描表都是 universal solenoid 的有限层因子观测，而全部相容有限层的逆极限恰恢复 \(\Sigma_{\mathbb P}\)。
\end{proposition}
\begin{proof}
由
\[
0\longrightarrow \ZZ\longrightarrow \QQ\longrightarrow \QQ/\ZZ\longrightarrow 0,
\qquad
\QQ/\ZZ\cong \bigoplus_{p\in\mathbb P}C_{p^\infty},
\]
对偶化即得所示短正合列。

再由
\[
\QQ=\varinjlim_{N\ge 1}(1/N)\ZZ
\]
（偏序由整除关系给出），Pontryagin 对偶把该直极限送为逆极限，从而
\[
\Sigma_{\mathbb P}
\cong
\varprojlim_{N\ge 1}\widehat{(1/N)\ZZ}.
\]
把 \(\widehat{(1/N)\ZZ}\) 识别为 \(\TT\) 后，若 \(N\mid M\)，则包含
\[
(1/N)\ZZ\hookrightarrow (1/M)\ZZ
\]
的对偶映射满足
\[
z\longmapsto z^{M/N}.
\]

第 \(1\) 条由定义立即得到：
\[
\Phi_q^{\cos}(x)+i\Phi_q^{\sin}(x)
=
\frac12\bigl(\chi_{1+q}(x)+\chi_{1-q}(x)\bigr)
+\frac12\bigl(\chi_{1+q}(x)-\chi_{1-q}(x)\bigr)
=
\chi_{1+q}(x).
\]

对第 \(2\) 条，由上一式可知 \(\Phi^{\mathrm{ph}}(x)\) 决定全部字符
\[
\chi_r(x)
\qquad (r\in\QQ),
\]
因为可取 \(r=1+q\)。而 \(\QQ\) 正是 \(\Sigma_{\mathbb P}\) 的对偶群，其字符分离点，故 \(\Phi^{\mathrm{ph}}\) 单射。连续性显然，而 \(\Sigma_{\mathbb P}\) 紧、\(\prod_{q\in\QQ}\CC^2\) Hausdorff，因此连续单射自动是到像上的同胚，即拓扑嵌入。

对第 \(3\) 条，若 \(q\in E\)，则 \(q\in H_{N_E}\)，从而
\[
1\pm q\in H_{N_E}.
\]
故 \(\chi_{1\pm q}\) 都只依赖于 \(\Sigma_{\mathbb P}\) 在 \(\widehat{H_{N_E}}\) 上的投影，也即只依赖于 \(\pi_{N_E}\)。于是 \(\Phi_q^{\cos}\) 与 \(\Phi_q^{\sin}\) 都经由 \(\pi_{N_E}\) 因子化，有限投影 \(\Phi_E^{\mathrm{ph}}\) 亦然。

最后，由 \(\QQ=\bigcup_{N\ge 1}(1/N)\ZZ\)，任一有理分母观测都落在某一有限层 \(H_N\) 中；而全部这些有限层通过整除偏序相容拼接，便恢复整个对偶群 \(\QQ\)，从而恢复 \(\Sigma_{\mathbb P}\)。证毕。
\end{proof}

\begin{corollary}[有理玫瑰的一阶字符承载对象由分母素数支撑唯一锁定]\label{cor:xi-time-part62d-rational-rose-minimal-prime-carrier}
\leanverified{paper\_xi\_time\_part62d\_rational\_rose\_minimal\_prime\_carrier}
设 \(q=n/d\in\QQ\) 为既约有理数，\(d\ge 1\)，并记 \(S(q)\subset\mathbb P\) 为 \(d\) 的素因子集合。则在定理 \ref{thm:xi-time-part62d-solenoid-kronecker-rose-lissajous-lift} 的一阶字符读出口径下，玫瑰轨道
\[
\gamma_q(\theta)=\cos(q\theta)e^{i\theta}
\]
可由 \(\Sigma_{S(q)}\) 实现，且不能由任何严格更小的 \(\Sigma_{S'}\)（\(S'\subsetneq S(q)\)）实现。并且
\[
p\in S(q)
\iff
\prod_{r\in S(q)}\ZZ_r\ \text{的}\ p\text{-主因子非平凡}.
\]
因此，最小承载对象的 prime-register 核完全恢复分母的素数支撑。
\end{corollary}
\begin{proof}
由定理 \ref{thm:xi-time-part62d-solenoid-kronecker-rose-lissajous-lift}，\(\gamma_q\) 的读出公式只使用字符 \(\chi_{1+q}\) 与 \(\chi_{1-q}\)。这两者在 \(\Sigma_S=\widehat{\ZZ[S^{-1}]}\) 上有定义，当且仅当
\[
1+q,\ 1-q\in \ZZ[S^{-1}],
\]
等价地，当且仅当 \(q\in \ZZ[S^{-1}]\)。而对既约有理数 \(q=n/d\)，这又等价于 \(d\) 的全部素因子都属于 \(S\)。故最小可行素数集恰为 \(S(q)\)。

若 \(S'\subsetneq S(q)\)，则 \(q\notin \ZZ[S'^{-1}]\)，从而上述一阶字符公式在 \(\Sigma_{S'}\) 上无法成立，因此 \(\Sigma_{S(q)}\) 在该读出口径下确为最小承载对象。

最后，由定理 \ref{thm:killo-localization-circle-dimension-prime-ledger-splitting}，\(\Sigma_{S(q)}\) 的核
\[
K_{S(q)}=\prod_{p\in S(q)}\ZZ_p
\]
满足
\[
p\in S(q)
\iff
K_{S(q)}\ \text{的}\ p\text{-主因子非平凡}.
\]
故分母的素数支撑可由 prime-register 核完全恢复。证毕。
\end{proof}

\begin{theorem}[有理玫瑰轨道的最小周期与花瓣数分类]\label{thm:xi-time-part62d-rational-rose-period-petal-classification}
\leanverified{paper\_xi\_time\_part62d\_rational\_rose\_period\_petal\_classification}
设 \(q=n/d\) 为正既约有理数，其中 \(n,d\in\NN\)。定义
\[
\gamma_q(\theta):=\cos\!\Bigl(\frac{n}{d}\theta\Bigr)e^{i\theta}.
\]
并把
\[
\gamma_q(\RR)\setminus\{0\}
\]
的连通分支闭包称为该轨道的花瓣。则其最小周期 \(T(q)\) 与花瓣数 \(P(q)\) 满足
\[
T(q)=
\begin{cases}
\pi d, & n\equiv d\equiv 1\pmod 2,\\[1mm]
2\pi d, & \text{否则},
\end{cases}
\qquad
P(q)=
\begin{cases}
n, & n\equiv d\equiv 1\pmod 2,\\[1mm]
2n, & \text{否则}.
\end{cases}
\]
\end{theorem}
\begin{proof}
把 \(\gamma_q\) 改写为
\[
\gamma_q(\theta)
=
\frac12\Bigl(
e^{i(1+n/d)\theta}+e^{i(1-n/d)\theta}
\Bigr).
\]
令
\[
\alpha:=d+n,
\qquad
\beta:=d-n.
\]
若 \(T>0\) 是其周期，则必须同时满足
\[
e^{i\alpha T/d}=1,
\qquad
e^{i\beta T/d}=1.
\]
故最小正周期为
\[
T(q)
=
\frac{2\pi d}{\gcd(\alpha,\beta)}
=
\frac{2\pi d}{\gcd(d+n,d-n)}
=
\frac{2\pi d}{\gcd(d+n,2n)}.
\]
由于 \((n,d)=1\)，有
\[
\gcd(d+n,2n)=\gcd(d+n,2).
\]
因此当且仅当 \(d+n\) 为偶数时分母为 \(2\)。又 \((n,d)=1\) 下，\(d+n\) 为偶数等价于 \(n,d\) 同为奇数。于是得到所示 \(T(q)\)。

现计数花瓣。原点恰在
\[
\theta_k:=\frac{d(2k+1)\pi}{2n}
\qquad (k\in\ZZ)
\]
处出现，因为此时 \(\cos(n\theta/d)=0\)。因此每个开区间
\[
I_k:=(\theta_k,\theta_{k+1})
\]
的像都落在 \(\gamma_q(\RR)\setminus\{0\}\) 的某个连通分支中，而全部花瓣恰由这些区间像的闭包给出。

再注意到步长 \(\pi d/n\) 的平移满足
\[
\gamma_q\!\Bigl(\theta+\frac{\pi d}{n}\Bigr)
=
\cos\!\Bigl(\frac{n}{d}\theta+\pi\Bigr)e^{i\theta}e^{i\pi d/n}
=
\omega\,\gamma_q(\theta),
\]
其中
\[
\omega:=-e^{i\pi d/n}=e^{i\pi(d+n)/n}.
\]
故相邻零区间上的像互为绕原点的旋转；全部花瓣正是一朵基准花瓣在循环群 \(\langle \omega\rangle\) 作用下的轨道。于是花瓣数等于 \(\omega\) 的阶。

而 \(\omega^m=1\) 当且仅当
\[
\frac{m(d+n)}{n}\in 2\ZZ.
\]
又因 \((n,d)=1\)，有 \(\gcd(d+n,n)=1\)。因此最小这样的 \(m\) 为
\[
m=
\begin{cases}
n, & d+n\ \text{为偶数},\\
2n, & d+n\ \text{为奇数}.
\end{cases}
\]
结合前述奇偶判别，即得
\[
P(q)=
\begin{cases}
n, & n\equiv d\equiv 1\pmod 2,\\
2n, & \text{否则}.
\end{cases}
\]
证毕。
\end{proof}

\begin{proposition}[Lissajous 自交点作为 Klein 四元反演像的轨道交会]\label{prop:xi-time-part62d-lissajous-klein-orbit-intersection}
\leanverified{paper\_xi\_time\_part62d\_lissajous\_klein\_orbit\_intersection}
设整数对 \((a,b)\neq (0,0)\)，相位 \(\delta\in\RR\)，并记
\[
L_{a,b,\delta}(t):=(\cos(at+\delta),\cos(bt)).
\]
定义其环面提升
\[
\widetilde L_{a,b,\delta}(t):=(e^{i(at+\delta)},e^{ibt})\in\TT^2,
\]
投影
\[
\Pi:\TT^2\to[-1,1]^2,
\qquad
\Pi(z_1,z_2):=(\Re z_1,\Re z_2),
\]
以及 Klein 四元群
\[
V_4:=\{(\varepsilon_1,\varepsilon_2):\ \varepsilon_j\in\{\pm 1\}\}
\]
在 \(\TT^2\) 上的作用
\[
(\varepsilon_1,\varepsilon_2)\cdot(z_1,z_2):=(z_1^{\varepsilon_1},z_2^{\varepsilon_2}).
\]
记
\[
g:=\gcd(|a|,|b|),
\qquad
T_{a,b}:=\frac{2\pi}{g}.
\]
则 \(L_{a,b,\delta}=\Pi\circ \widetilde L_{a,b,\delta}\)，且对任意
\[
t,t'\in[0,T_{a,b}),
\qquad
t\neq t',
\]
有严格等价
\[
L_{a,b,\delta}(t)=L_{a,b,\delta}(t')
\iff
\exists\,(\varepsilon_1,\varepsilon_2)\in V_4\setminus\{(1,1)\}
\ \text{使}\ 
\widetilde L_{a,b,\delta}(t)
=
(\varepsilon_1,\varepsilon_2)\cdot
\widetilde L_{a,b,\delta}(t').
\]
因此，Lissajous 的全部自交点都来源于提升轨道与其三条非平凡反演像之间的交会。
\end{proposition}
\begin{proof}
恒等式 \(L_{a,b,\delta}=\Pi\circ \widetilde L_{a,b,\delta}\) 由 \(\Re(e^{i\theta})=\cos\theta\) 立即得到。

现取 \(t,t'\in[0,T_{a,b})\)。若
\[
L_{a,b,\delta}(t)=L_{a,b,\delta}(t'),
\]
则
\[
\Re\bigl(e^{i(at+\delta)}\bigr)=\Re\bigl(e^{i(at'+\delta)}\bigr),
\qquad
\Re\bigl(e^{ibt}\bigr)=\Re\bigl(e^{ibt'}\bigr).
\]
在单位圆上，两个点实部相同当且仅当二者相等或互为逆元，因此存在
\[
\varepsilon_1,\varepsilon_2\in\{\pm1\}
\]
使
\[
e^{i(at+\delta)}=\bigl(e^{i(at'+\delta)}\bigr)^{\varepsilon_1},
\qquad
e^{ibt}=\bigl(e^{ibt'}\bigr)^{\varepsilon_2}.
\]
这正是
\[
\widetilde L_{a,b,\delta}(t)
=
(\varepsilon_1,\varepsilon_2)\cdot
\widetilde L_{a,b,\delta}(t').
\]

若 \((\varepsilon_1,\varepsilon_2)=(1,1)\)，则
\[
 a(t-t')\in 2\pi\ZZ,
\qquad
 b(t-t')\in 2\pi\ZZ.
\]
取 Bézout 系数 \(u,v\in\ZZ\) 使
\[
ua+vb=g.
\]
则
\[
g(t-t')=u\,a(t-t')+v\,b(t-t')\in 2\pi\ZZ,
\]
从而
\[
t-t'\in (2\pi/g)\ZZ=T_{a,b}\ZZ.
\]
由于 \(t,t'\in[0,T_{a,b})\)，只能有 \(t=t'\)，与假设矛盾。故此时 \((\varepsilon_1,\varepsilon_2)\neq (1,1)\)。

反向推出同样直接：若存在非平凡 \((\varepsilon_1,\varepsilon_2)\) 使
\[
\widetilde L_{a,b,\delta}(t)
=
(\varepsilon_1,\varepsilon_2)\cdot
\widetilde L_{a,b,\delta}(t'),
\]
则两坐标分别具有相同实部，因而
\[
\Pi\bigl(\widetilde L_{a,b,\delta}(t)\bigr)
=
\Pi\bigl(\widetilde L_{a,b,\delta}(t')\bigr),
\]
即 \(L_{a,b,\delta}(t)=L_{a,b,\delta}(t')\)。证毕。
\end{proof}

\begin{theorem}[Fibonacci 挠点并集饱和单位圆的全部有限挠点]\label{thm:xi-time-part62da-fibonacci-torsion-saturates-circle}
\leanverified{paper\_xi\_time\_part62da\_fibonacci\_torsion\_saturates\_circle}
定义 Fibonacci 挠点并集
\[
\mathrm{Tor}_F(\TT):=\bigcup_{m\ge 1}\mu_{F_{m+2}},
\]
其中 \(\mu_n\subset\TT\) 表示 \(n\) 次单位根群。则有严格恒等式
\[
\boxed{\;
\mathrm{Tor}_F(\TT)=\mathrm{Tor}(\TT):=\bigcup_{N\ge 1}\mu_N.\;}
\]
\end{theorem}
\begin{proof}
显然
\[
\mathrm{Tor}_F(\TT)\subseteq \mathrm{Tor}(\TT).
\]
反向包含方面，任取 \(\zeta\in\mu_N\)。由定理 \ref{thm:xi-visible-arithmetic-fibonacci-adic-profinite-coincidence} 的证明，对该 \(N\) 存在某个 \(m\ge 1\) 使
\[
N\mid F_{m+2}.
\]
于是
\[
\zeta^{F_{m+2}}=1,
\]
故 \(\zeta\in \mu_{F_{m+2}}\subseteq \mathrm{Tor}_F(\TT)\)。从而 \(\mathrm{Tor}(\TT)\subseteq \mathrm{Tor}_F(\TT)\)，两侧相等。证毕。
\end{proof}

\begin{theorem}[Fibonacci 分母子系统恢复 universal solenoid]\label{thm:xi-time-part62da-fibonacci-solenoid-cofinal-universal}
\leanverified{paper\_xi\_time\_part62da\_fibonacci\_solenoid\_cofinal\_universal}
记
\[
S_F:=\{F_{m+2}:\ m\ge 1\}\subset \NN.
\]
把 \(S_F\) 视为整除偏序 \((\NN,\mid)\) 的共终子系统，并定义对应逆极限
\[
\Sigma_F:=\varprojlim_{N\in S_F}\TT,
\]
其中当 \(N\mid N'\) 时连接映射取为 \(z\mapsto z^{N'/N}\)。则有规范同构
\[
\boxed{\;
\Sigma_F\cong \Sigma_{\mathbb P}=\widehat{\QQ}\cong \varprojlim_{N\ge 1}\TT.\;}
\]
\end{theorem}
\begin{proof}
由定理 \ref{thm:xi-visible-arithmetic-fibonacci-adic-profinite-coincidence} 的证明，对任意 \(N\ge 1\) 都存在某个 \(m\ge 1\) 使
\[
N\mid F_{m+2}.
\]
故 \(S_F\) 在整除偏序中对 \(\NN\) 共终。

另一方面，由命题 \ref{prop:xi-time-part62d-universal-solenoid-phase-completed-scan}，
\[
\Sigma_{\mathbb P}\cong \varprojlim_{N\ge 1}\TT
\]
且连接映射正是整除关系诱导的幂映射。逆极限对共终限制不变，故
\[
\varprojlim_{N\ge 1}\TT\cong \varprojlim_{N\in S_F}\TT=\Sigma_F.
\]
合并即得所示规范同构。证毕。
\end{proof}

\begin{theorem}[Fibonacci 分母子系统的单位群逆极限恢复 profinite 单位群]\label{thm:xi-time-part62da-fibonacci-unit-limit-zhat-units}
定义
\[
U_F:=\varprojlim_{m}(\ZZ/F_{m+2}\ZZ)^\times,
\]
其中连接映射由整除关系下的标准降模同态给出。则存在规范拓扑群同构
\[
\boxed{\;
U_F\cong \widehat{\ZZ}^{\,\times}.\;}
\]
\end{theorem}
\begin{proof}
由定理 \ref{thm:xi-time-part62da-fibonacci-solenoid-cofinal-universal}，集合 \(S_F=\{F_{m+2}\}\) 在整除偏序中共终。故对单位群系统同样可作共终限制：
\[
\varprojlim_{N\ge 1}(\ZZ/N\ZZ)^\times
\cong
\varprojlim_{N\in S_F}(\ZZ/N\ZZ)^\times
=
U_F.
\]
而左端正是 \(\widehat{\ZZ}^{\,\times}\) 的经典表示。证毕。
\end{proof}

\begin{corollary}[Dirichlet 扭曲在 Fibonacci 分辨率上内禀实现]\label{cor:xi-time-part62da-dirichlet-twists-intrinsic-on-fibonacci-levels}
\leanverified{paper\_xi\_time\_part62da\_dirichlet\_twists\_intrinsic\_on\_fibonacci\_levels}
任取正整数 \(N\) 与 Dirichlet 特征
\[
\chi:(\ZZ/N\ZZ)^\times\to \CC^\times.
\]
则存在某个 \(m\ge 1\) 与群同态
\[
\widetilde\chi:(\ZZ/F_{m+2}\ZZ)^\times\to \CC^\times
\]
使得
\[
\widetilde\chi=\chi\circ \rho_{m,N},
\]
其中 \(\rho_{m,N}\) 为标准降模映射
\[
(\ZZ/F_{m+2}\ZZ)^\times\twoheadrightarrow (\ZZ/N\ZZ)^\times.
\]
因此，任意有限模 Dirichlet 扭曲都可在某一级 Fibonacci 分辨率上内禀实现；命题 \ref{prop:real-input-40-output-dirichlet-nonrh} 中的单位根扭曲只是其中最基本的特例。
\end{corollary}
\begin{proof}
由定理 \ref{thm:xi-visible-arithmetic-fibonacci-adic-profinite-coincidence} 的证明，可取 \(m\ge 1\) 使 \(N\mid F_{m+2}\)。于是标准降模映射
\[
\rho_{m,N}:(\ZZ/F_{m+2}\ZZ)^\times\twoheadrightarrow (\ZZ/N\ZZ)^\times
\]
为满射。令
\[
\widetilde\chi:=\chi\circ \rho_{m,N}
\]
便得到在 Fibonacci 模数层上的提升。再由定理 \ref{thm:fold-congruence-mod-semirings}，
\[
\Omega_m/\!\equiv_m\cong \ZZ/(F_{m+2}\ZZ),
\]
故该扭曲首先在可见同余层的单位扇区上给出规范乘法权重；再经由商映射
\[
\Omega_m\twoheadrightarrow \Omega_m/\!\equiv_m
\]
即可拉回到对应的微观态单位扇区。若需把它扩展为定义在全部微观态上的权重，只需再对非单位扇区作任意固定延拓。证毕。
\end{proof}

\begin{proposition}[有理玫瑰在原点的唯一聚合奇异性与切锥方向数]\label{prop:xi-time-part62da-rational-rose-origin-tangent-cone}
\leanverified{paper\_xi\_time\_part62da\_rational\_rose\_origin\_tangent\_cone}
设 \(n,d\in\NN\) 互素，且 \(d>n\)。定义复平面参数曲线
\[
z_{n,d}(t):=e^{idt}\cos(nt)
\qquad (t\in\RR).
\]
则：
\begin{enumerate}
  \item 参数曲线 \(z_{n,d}\) 处处浸入，但其像在原点 \(0\) 处存在自交聚合；并且原点是唯一由 \(\cos(nt)=0\) 所产生的聚合像点。
  \item 原点处不同切线方向的条数恰为
  \[
  \boxed{\;
  \#\mathrm{Tan}_0(z_{n,d})
  =
  n.\;}
  \]
  全部方向在 \([0,\pi)\) 内等角分布。
\end{enumerate}
\end{proposition}
\begin{proof}
首先，\(z_{n,d}(t)=0\) 当且仅当 \(\cos(nt)=0\)，即
\[
t=t_k:=\frac{(2k+1)\pi}{2n}
\qquad (k\in\ZZ).
\]
因此原点恰由这一离散时刻族聚合得到。另一方面，若 \(z_{n,d}(t)=0\)，则必满足 \(\cos(nt)=0\)，故在该参数表示下原点是唯一零像聚合点。

再求导得
\[
z'_{n,d}(t)
=
e^{idt}\bigl(-n\sin(nt)+id\cos(nt)\bigr).
\]
在 \(t=t_k\) 处有 \(\cos(nt_k)=0\) 且 \(\sin(nt_k)=\pm 1\)，从而
\[
z'_{n,d}(t_k)=\mp n\,e^{idt_k}\neq 0.
\]
故曲线处处浸入，原点处的奇异性来自多分支聚合，而非尖点。

原点处分支的切线方向由
\[
\arg z'_{n,d}(t_k)\equiv dt_k\pmod\pi
\]
决定。于是不同方向数等于集合
\[
\left\{\frac{d(2k+1)\pi}{2n}\bmod \pi:\ 0\le k<2n\right\}
\]
的基数。两项给出同一方向当且仅当
\[
d(2k+1)\equiv d(2k'+1)\pmod{2n},
\]
即
\[
n\mid d(k-k').
\]
由于 \(\gcd(n,d)=1\)，这等价于
\[
n\mid (k-k').
\]
因此在 \(0\le k<2n\) 的范围内，每个方向类恰由 \(k\) 与 \(k+n\) 这两个参数点组成，故不同切线方向总数为 \(n\)。

再看其分布。把方向写成
\[
\frac{d(2k+1)\pi}{2n}
\equiv
\frac{d\pi}{2n}+\frac{dk\pi}{n}
\pmod\pi.
\]
由于 \(\gcd(d,n)=1\)，乘以 \(d\) 在 \(\ZZ/n\ZZ\) 上只是一个置换，因此这些方向与
\[
\frac{\pi}{2n}+\frac{k\pi}{n}
\qquad (0\le k<n)
\]
给出的 \(n\) 个方向完全相同，从而在 \([0,\pi)\) 内等角分布。证毕。
\end{proof}

\leanverified{paper\_xi\_time\_part62da\_lissajous\_two\_nonimmersed\_points}
\begin{proposition}[Lissajous 的非沉浸点判别与精确二点数目]\label{prop:xi-time-part62da-lissajous-two-nonimmersed-points}
设 \(a,b\in\NN\) 互素，\(\delta\in\RR\)，并定义
\[
\gamma_{a,b,\delta}(t):=(\cos(at+\delta),\cos(bt))
\qquad (t\in\RR/2\pi\ZZ).
\]
则 \(\gamma_{a,b,\delta}\) 存在非沉浸点当且仅当
\[
\boxed{\;
\delta\in \frac{\pi}{b}\ZZ.\;}
\]
在此情形下，非沉浸点恰有 \(2\) 个。更精确地，若 \(k\) 是模 \(b\) 下满足
\[
ak+\frac{b\delta}{\pi}\equiv 0\pmod b
\]
的唯一剩余类，则对应参数值为
\[
\boxed{\;
t=\frac{k\pi}{b},
\qquad
t=\frac{(k+b)\pi}{b}\pmod{2\pi}.\;}
\]
\end{proposition}
\begin{proof}
由定理 \ref{thm:xi-lissajous-phase-singular-ring-immersion}，
\[
\gamma'_{a,b,\delta}(t)=0
\]
当且仅当
\[
\sin(bt)=0,
\qquad
\sin(at+\delta)=0.
\]
前者给出
\[
t=\frac{k\pi}{b}
\qquad(k\in\ZZ).
\]
代入第二式，得到同余条件
\[
ak+\frac{b\delta}{\pi}\equiv 0\pmod b.
\]
由于 \(\gcd(a,b)=1\)，\(a\) 在模 \(b\) 下可逆，因此该同余有解当且仅当
\[
\frac{b\delta}{\pi}\in \ZZ,
\]
即 \(\delta\in \frac{\pi}{b}\ZZ\)；并且解 \(k\) 在模 \(b\) 下唯一。

在 \([0,2\pi)\) 中，与这一模 \(b\) 解对应的参数恰为
\[
\frac{k\pi}{b}
\quad\text{与}\quad
\frac{(k+b)\pi}{b},
\]
二者相差 \(\pi\)，给出两个不同点。除此之外不再有别的 \(k\) 模 \(b\) 解，故非沉浸点恰有两处。证毕。
\end{proof}


在上述命题与命题 \ref{prop:xi-time-part62d-universal-solenoid-phase-completed-scan} 之间，实部通道给出经典玫瑰/Lissajous 的几何阴影，而与之互补的虚部通道则正是相位补全所必需的附加数据。由此可把 RH 一侧的圆盘完成化接口组织为如下猜想性口径。

\begin{conjecture}[虚部读出下的 prime-register 核消隐与 RH 的 Toeplitz 因子化]\label{conj:xi-time-part62d-imaginary-channel-prime-kernel-rh}
沿用命题 \ref{prop:xi-time-part62d-universal-solenoid-phase-completed-scan} 的相位补全坐标与定理 \ref{thm:xi-disk-zero-pole-index-reduction} 的圆盘完成化接口。设
\[
\Sigma_{\mathbb P}=\widehat{\QQ}
\twoheadrightarrow
\TT
\]
为 universal solenoid 的圆商。则 RH 应等价于如下核消隐原则：在与 \(\mathscr C_\zeta\) 对应的 Weyl--Carath\'eodory/Toeplitz 可见层上，全部边界数据都经由该圆商因子化；等价地，凡对 prime-register 核
\[
\prod_{p\in\mathbb P}\ZZ_p
\]
真正敏感的虚部通道分量都必须消失。若某一有限层仍残留此类核敏感分量，则相应 Toeplitz 截断不再维持 Carath\'eodory--PSD 正性，并应通过定理 \ref{thm:xi-disk-zero-pole-index-reduction} 诱发盘内零点，从而否定 RH。
\end{conjecture}

\endinput
